%
\documentclass[a4paper]{article}
\usepackage{Sweave}
\usepackage[utf8]{inputenc} %encodage du fichier source
\usepackage[T1]{fontenc}  %gestion des accents (pour les pdf)
\usepackage[french]{babel}
\usepackage{float}
\usepackage[document]{ragged2e} % pour que le texte soit aligné à gauche
\usepackage[section]{placeins}%The placeins package provides the command \FloatBarrier
\usepackage{amsmath} % equation numbering
\usepackage[colorlinks=true,urlcolor=blue,linkcolor=black,citecolor=black]{hyperref}
\usepackage{wrapfig}
\usepackage{subcaption} % pour les subfigures
\usepackage{vmargin}
\setpapersize{A4}
\setmarginsrb{15mm}{15mm}{15mm}{15mm}{12pt}{0mm}{0pt}{11mm}
\pagestyle{empty}

\title{Sortie des parametres Openbugs Modele
2026\_07 (Mise à jour des données 2025)}
\author{Marion Legrand \& Etienne Prévost}

\begin{document}
\input{SortieOpenbugs_parameters-concordance}


\maketitle


\clearpage
\section{Survie du juvénile à l'adulte}

%========================
\subsection{s\_juv2ad}
%=================================================================================

\begin{figure}[ht]
     \includegraphics{D:/Documents/Workspace_eclipse/ModeleDynamiquePop/img/2026_07/Model_influence_Q/s_juv2ad.png}%{a_juv.png}
     \caption{s\_juv2ad : la mediane est representée}
     \label{s_juv2ad}
\end{figure}

s\_juv2ad représente le taux de transition entre le juvénile 0+ et l'adulte.

\clearpage

%=================================================================
\subsection{Influence des débits sur la survie juvénile-adulte}
%=================================================================
%========================
\subsubsection{m\_surv}
%=================================================================================

\begin{figure}[ht]
     \includegraphics{D:/Documents/Workspace_eclipse/ModeleDynamiquePop/img/2026_07/Model_influence_Q/m_surv.png}%{a_juv.png}
     \caption{m\_surv}
     \label{m_surv}
\end{figure}

\input{D:/Documents/Workspace_eclipse/ModeleDynamiquePop/tab/2026_07/Model_influence_Q/m_surv.tex}
m\_surv représente la survie moyenne. Il est utilisé dans la relation qui lie la survie du juvénile à l'adulte aux débits.

\clearpage

%=================================================================================
\subsubsection{beta\_Q\_VB\_deval}
%=================================================================================

\begin{figure}[ht]
     \includegraphics{D:/Documents/Workspace_eclipse/ModeleDynamiquePop/img/2026_07/Model_influence_Q/beta_Q_VB_deval.png}
     \caption{beta\_Q\_VB\_deval}
     \label{beta_Q_VB_deval}
\end{figure}

\input{D:/Documents/Workspace_eclipse/ModeleDynamiquePop/tab/2026_07/Model_influence_Q/beta_Q_VB_deval.tex}

beta\_Q\_VB\_deval représente le coefficient directeur de la droite modélisant l'influence des débits sur la survie. Les débits sont pris à Vieille-Brioude lors de la dévalaison (mars-avril).

\clearpage

%=================================================================================
\subsubsection{beta\_Q\_MSL\_mont}
%=================================================================================

\begin{figure}[ht]
     \includegraphics{D:/Documents/Workspace_eclipse/ModeleDynamiquePop/img/2026_07/Model_influence_Q/beta_Q_MSL_mont.png}
     \caption{beta\_Q\_MSL\_mont}
     \label{beta_Q_MSL_mont}
\end{figure}

\input{D:/Documents/Workspace_eclipse/ModeleDynamiquePop/tab/2026_07/Model_influence_Q/beta_Q_MSL_mont.tex}

beta\_Q\_MSL\_mont représente le coefficient directeur de la droite modélisant l'influence des débits sur la survie. Les débits sont pris à Montjean-sur-Loire lors de la montaison (du 15/12 au 15/03).

\clearpage

%=================================================================================
\subsubsection{beta\_Q\_MSL\_deval}
%=================================================================================

\begin{figure}[ht]
     \includegraphics{D:/Documents/Workspace_eclipse/ModeleDynamiquePop/img/2026_07/Model_influence_Q/beta_Q_MSL_deval.png}
     \caption{beta\_Q\_MSL\_deval}
     \label{beta_Q_MSL_deval}
\end{figure}

\input{D:/Documents/Workspace_eclipse/ModeleDynamiquePop/tab/2026_07/Model_influence_Q/beta_Q_MSL_deval.tex}

beta\_Q\_MSL\_deval représente le coefficient directeur de la droite modélisant l'influence des débits sur la survie. Les débits sont pris à Montjean-sur-Loire lors de la dévalaison (du 15/03 au 15/05).

\clearpage


%=================================================================================
\subsubsection{corrplot\_paramètre\_débit}
%=================================================================================



\begin{figure}[ht]
     \includegraphics{D:/Documents/Workspace_eclipse/ModeleDynamiquePop/img/2026_07/Model_influence_Q/corr_Q_surv.png}%{a_juv.png}
    % \caption{corr\_Q\_survie : débits utilisés dans le modèle linéaire expliquant le taux de transition 0\ensuremath{+}/adulte par les débits (Montjean-sur Loire à la dévalaison, Montjean-sur Loire à la montaison et Vieille-Brioude à la dévalaison)}
     \label{corr_Q_surv}
\end{figure}


\clearpage


\subsubsection{Représentation des séries de débits utilisées dans le modèle}


\begin{figure}[ht]
     \includegraphics{D:/Documents/Workspace_eclipse/ModeleDynamiquePop/img/2026_07/Model_influence_Q/plot_Q_surv.png}
     %\caption{représentation des séries de débits utilisées dans le modèle pour
    % expliquer le taux de transition 0\ensuremath{+}/adulte}
     \label{serie_Q_surv}
\end{figure}

\clearpage

\subsubsection{corrplot\_série\_débit}

\begin{figure}[ht]
     \includegraphics{D:/Documents/Workspace_eclipse/ModeleDynamiquePop/img/2026_07/Model_influence_Q/corr_Q_surv_data.png}%{a_juv.png}
     %\caption{Débits utilisés dans le modèle linéaire expliquant le taux de
     %transition 0\ensuremath{+}/adulte par les débits
     %(Montjean-sur-Loire montaison et dévalaison et Vieille-Brioude dévalaison)}
     \label{corr_Q_rep_surv}
\end{figure}


\clearpage

%============================================
%\subsection{Paramètres du modèle de mélange}
%============================================
%=================================================================================
%\subsubsection{I\_mm}
%================================================================================

% \begin{figure}[ht]
%      \includegraphics{D:/Documents/Workspace_eclipse/ModeleDynamiquePop/img/2026_07/Model_influence_Q/I_mm.png}%{a_juv.png}
%      \caption{I\_mm : représentation des moyennes}
%      \label{I_mm}
% \end{figure}
%
% I\_mm représente l'indicatrice des bonnes et des mauvaises années (1=mauvaise année / 0= bonne année) pour les retours d'adultes à Vichy.
%
% \clearpage
%=================================================================================
%\subsubsection{ratio\_surv}
%================================================================================

% \begin{figure}[ht]
%      \includegraphics{D:/Documents/Workspace_eclipse/ModeleDynamiquePop/img/2026_07/Model_influence_Q/ratio_surv.png}
%      \caption{ratio\_surv}
%      \label{ratio_surv}
% \end{figure}
%
% \input{D:/Documents/Workspace_eclipse/ModeleDynamiquePop/tab/2026_07/Model_influence_Q/ratio_surv.tex}
% \clearpage


%=================================================================================
%\subsubsection{p\_surv}
%================================================================================
%
% \begin{figure}[ht]
%      \includegraphics{D:/Documents/Workspace_eclipse/ModeleDynamiquePop/img/2026_07/Model_influence_Q/p_surv.png}
%      \caption{p\_surv}
%      \label{p_surv}
% \end{figure}
%
% \input{D:/Documents/Workspace_eclipse/ModeleDynamiquePop/tab/2026_07/Model_influence_Q/p_surv.tex}
%
% p\_surv représente le paramètre de surmortalité lorsque l'on se trouve dans une mauvaise année.
%
% \clearpage

%=================================================================================
\subsection{sigma\_vichy}
%=================================================================================

\begin{figure}[ht]
     \includegraphics{D:/Documents/Workspace_eclipse/ModeleDynamiquePop/img/2026_07/Model_influence_Q/sigma_vichy.png}%{a_juv.png}
     \caption{sigma\_vichy}
     \label{sigma_vichy}
\end{figure}

\input{D:/Documents/Workspace_eclipse/ModeleDynamiquePop/tab/2026_07/Model_influence_Q/sigma_vichy.tex}

sigma\_vichy représente l'écart-type des estimations d'adultes de retour à Vichy.
\clearpage


%====================================
\subsection{N\_Vichy}
%====================================


\begin{figure}[ht]
     \includegraphics{D:/Documents/Workspace_eclipse/ModeleDynamiquePop/img/2026_07/Model_influence_Q/N_vichy.png}%{a_juv.png}
     \caption{N\_vichy}
     \label{N_vichy}
\end{figure}

N\_Vichy représente les adultes de retour à Vichy. Lorsque la station de comptage est en fonctionnement, la donnée observée est représentée sous forme de point. Sinon, les estimations des effectifs sont représentées sous forme de boxplot.

\clearpage




%=================================================================================
\subsection{res\_vichy\_standardise}
%================================================================================

\begin{figure}[ht]
     \includegraphics{D:/Documents/Workspace_eclipse/ModeleDynamiquePop/img/2026_07/Model_influence_Q/res_vichy_std.png}%{a_juv.png}
     \caption{res\_vichy\_std}
     \label{res_vichy_std}
\end{figure}

res\_vichy\_std représente les résidus standardisés des retours d'adultes à Vichy.

\clearpage





\section{Répartition des adultes}




%====================================
\subsection{N\_Alagnon}
%====================================

\begin{figure}[ht]
     \includegraphics{D:/Documents/Workspace_eclipse/ModeleDynamiquePop/img/2026_07/Model_influence_Q/N_alagnon.png}%{a_juv.png}
     \caption{N\_alagnon}
     \label{N_alagnon}
\end{figure}

N\_Alagnon représente les adultes migrant sur le cours d'eau Alagnon.

\clearpage

%====================================
\subsection{N\_Langeac}
%====================================

\begin{figure}[ht]
     \includegraphics{D:/Documents/Workspace_eclipse/ModeleDynamiquePop/img/2026_07/Model_influence_Q/N_langeac.png}%{a_juv.png}
     \caption{N\_langeac}
     \label{N_langeac}
\end{figure}

N\_Langeac représente les effectifs d'adultes ayant franchi Langeac. Ils sont calculés en tirant dans une binomiale avec probabilité de franchir Langeac et effectif égal à ce qui est passé à Vichy moins les captures pour la pisciculture moins les poissons qui ont migré sur l'Alagnon. La station de Langeac est considérée comme partielle dans le modèle. Ainsi, même lorsque la station de comptage est en fonctionnement, les effectifs sont estimés. A titre d'indication, les données observées à Langeac sont représentées sous forme de point.

\clearpage

%====================================
\subsection{N\_Poutès}
%====================================

\begin{figure}[ht]
     \includegraphics{D:/Documents/Workspace_eclipse/ModeleDynamiquePop/img/2026_07/Model_influence_Q/N_poutes.png}%{a_juv.png}
     \caption{N\_poutes}
     \label{N_poutes}
\end{figure}

N\_Poutès représente les effectifs de saumon ayant franchi Poutès. Durant la période où la station de comptage a fonctionné, les effectifs comptés sont exhaustifs et représentés par des points noirs. Depuis le démarrage des travaux (2019) et le réaménagement de l'ouvrage (2021), la période de transparence ne permet plus le comptage. Les effectifs sont estimés (boxplot).

\clearpage


%=================================================================================
\subsection{Adultes restant sur les différents secteurs de l'Allier}
%=================================================================================
Les adultes restant sur l'Alagnon correspondent au paramètre N\_alagnon (voir plus haut), et ceux restant en amont de Poutès correspondent au paramètre N\_poutes (voir plus haut). En revanche, pour le secteur Vichy-Langeac et Langeac-Poutès, les effectifs restant sur ces secteurs peuvent être estimés.

%=================================================================================
\subsubsection{Adultes restant sur Vichy-Langeac}
%=================================================================================



\begin{figure}[ht]
     \includegraphics{D:/Documents/Workspace_eclipse/ModeleDynamiquePop/img/2026_07/Model_influence_Q/adultes_vichy.png}%{a_juv.png}
     \caption{Adultes restant sur le secteur Vichy-Langeac}
     \label{adultes_vichy_langeac}
\end{figure}

\clearpage

%=================================================================================
\subsubsection{Adultes restant sur Langeac-Poutès}
%=================================================================================


\begin{figure}[ht]
     \includegraphics{D:/Documents/Workspace_eclipse/ModeleDynamiquePop/img/2026_07/Model_influence_Q/adultes_langeac.png}%{a_juv.png}
     \caption{Adultes restant sur le secteur Langeac-Poutès}
     \label{adultes_langeac_poutes}
\end{figure}

Le nombre d'adultes restant sur le secteur Langeac-Poutès est estimé à partir des effectifs qui sont passés à Vichy mais ne sont pas allés sur l'Alagnon et n'ont pas été capturés par la pisciculture et qui n'ont pas non plus franchi Poutès (on soustrait les effectifs comptés ou estimés à Poutès). Lorsque les stations de Langeac et Poutès ont fonctionné sur la même période on soustrait les effectifs comptés à Poutès à ceux de Langeac (pour rappel, la station de Langeac est considérée comme une station non-exhaustive dans le modèle).

\clearpage

%=================================================================================
\subsubsection{Pourcentage d'adultes dans les différents secteurs}
%=================================================================================


\begin{figure}[ht]
     \includegraphics{D:/Documents/Workspace_eclipse/ModeleDynamiquePop/img/2026_07/Model_influence_Q/proportion_adultes.png}%{a_juv.png}
     \caption{proportion des adultes dans les différents secteurs}
     \label{proportion_adultes}
\end{figure}


Les proportions dans les différents secteurs, sont calculés à partir des saumons estimés (avant 1997) ou comptés à Vichy, ainsi que des adultes estimés dans les différents secteurs du modèle.

\clearpage
%=================================================================================
\subsection{p\_comptage\_langeac}
%=================================================================================

\begin{figure}[ht]
     \includegraphics{D:/Documents/Workspace_eclipse/ModeleDynamiquePop/img/2026_07/Model_influence_Q/p_comptage_L.png}%{a_juv.png}
     \caption{p\_comptage\_L : les boxplots verts représentent les probabilités
     d'être comptés pour les années où la station de Langeac est en
     fonctionnement. En dehors de ces années, les boxplots sont gris}
     \label{p_comptage_L}
\end{figure}


p\_comptage\_L représente la probabilité qu'a un poisson franchissant
Langeac d'être compté

\clearpage

%=================================================================================
\subsection{N\_langeac\_obs}
%=================================================================================

\begin{figure}[ht]
     \includegraphics{D:/Documents/Workspace_eclipse/ModeleDynamiquePop/img/2026_07/Model_influence_Q/N_langeac_obs.png}%{a_juv.png}
     \caption{N\_langeac\_obs : les points représentent les effectifs comptés à
     Langeac (station de comptage en service). Les boxplots sont les
     estimations des effectifs comptés à Langeac lorsque la station n'est pas
     en service}
     \label{N_langeac_obs}
\end{figure}


N\_langeac\_obs représente les estimations des effectifs qui auraient été
comptés à Langeac si la station avait été en service

\clearpage
%=================================================================================
\subsection{rho\_station}
%=================================================================================

\begin{figure}[ht]
     \includegraphics{D:/Documents/Workspace_eclipse/ModeleDynamiquePop/img/2026_07/Model_influence_Q/rho_station.png}%{a_juv.png}
     \caption{rho\_station}
     \label{rho_station}
\end{figure}

\input{D:/Documents/Workspace_eclipse/ModeleDynamiquePop/tab/2026_07/Model_influence_Q/rho_station.tex}

rho\_station représente le facteur de pondération entre la répartition libre idéale en fonction des habitats disponibles et le homing. Si la médiane est supérieure à 0.50 cela indique que la présence de surfaces productives en amont joue un rôle prépondérant par rapport au homing.

\clearpage

%=================================================================================
\subsection{adjust\_p\_A}
%=================================================================================

\begin{figure}[ht]
     \includegraphics{D:/Documents/Workspace_eclipse/ModeleDynamiquePop/img/2026_07/Model_influence_Q/adjust_p_A.png}%{a_juv.png}
     \caption{adjust\_p\_A}
     \label{adjust_p_A}
\end{figure}

\input{D:/Documents/Workspace_eclipse/ModeleDynamiquePop/tab/2026_07/Model_influence_Q/adjust_p_A.tex}

adjust\_p\_A représente un paramètre de blocage venant contrarier la répartition idéale libre en fonction des habitats disponibles et du homing pour le secteur Alagnon.

\clearpage

%=================================================================================
\subsection{adjust\_p\_L}
%=================================================================================

\begin{figure}[ht]
     \includegraphics{D:/Documents/Workspace_eclipse/ModeleDynamiquePop/img/2026_07/Model_influence_Q/adjust_p_L.png}%{a_juv.png}
     \caption{adjust\_p\_L}
     \label{adjust_p_L}
\end{figure}

\input{D:/Documents/Workspace_eclipse/ModeleDynamiquePop/tab/2026_07/Model_influence_Q/adjust_p_L.tex}

adjust\_p\_L représente un paramètre de blocage venant contrarier la répartition idéale libre en fonction des habitats disponibles et du homing pour le secteur Langeac-Poutès.

\clearpage

%=================================================================================
\subsection{adjust\_p\_P}
%=================================================================================

\begin{figure}[ht]
     \includegraphics{D:/Documents/Workspace_eclipse/ModeleDynamiquePop/img/2026_07/Model_influence_Q/adjust_p_P.png}%{a_juv.png}
     \caption{adjust\_p\_P}
     \label{adjust_p_P}
\end{figure}

\input{D:/Documents/Workspace_eclipse/ModeleDynamiquePop/tab/2026_07/Model_influence_Q/adjust_p_P.tex}

adjust\_p\_P représente un paramètre de blocage venant contrarier la répartition idéale libre en fonction des habitats disponibles et du homing pour le secteur amont de Poutès.

\clearpage

%=================================================================================
\subsection{res\_p\_alagnon\_standardisé}
%================================================================================

\begin{figure}[ht]
     \includegraphics{D:/Documents/Workspace_eclipse/ModeleDynamiquePop/img/2026_07/Model_influence_Q/res_p_alagnon_std.png}%{a_juv.png}
     \caption{res\_p\_alagnon\_std}
     \label{res_p_alagnon_std}
\end{figure}

res\_p\_alagnon\_standardisé représente les résidus standardisés du paramètre de blocage de la migration pour le secteur Alagnon.

\clearpage



%=================================================================================
\subsection{res\_p\_langeac\_standardisé}
%================================================================================

\begin{figure}[ht]
     \includegraphics{D:/Documents/Workspace_eclipse/ModeleDynamiquePop/img/2026_07/Model_influence_Q/res_p_langeac_std.png}%{a_juv.png}
     \caption{res\_p\_langeac\_std}
     \label{res_p_langeac_std}
\end{figure}

res\_p\_langeac\_standardisé représente les résidus standardisés du paramètre de blocage de la migration pour le secteur Langeac-Poutès.

\clearpage


%=================================================================================
\subsection{res\_p\_poutes\_standardisé}
%================================================================================

\begin{figure}[ht]
     \includegraphics{D:/Documents/Workspace_eclipse/ModeleDynamiquePop/img/2026_07/Model_influence_Q/res_p_poutes_std.png}%{a_juv.png}
     \caption{res\_p\_poutes\_std}
     \label{res_p_poutes_std}
\end{figure}

res\_p\_poutes\_standardisé représente les résidus standardisés du paramètre de blocage de la migration pour le secteur amont de Poutès.

\clearpage


%=================================================================================
\subsection{sigma\_p\_alagnon}
%=================================================================================

\begin{figure}[ht]
     \includegraphics{D:/Documents/Workspace_eclipse/ModeleDynamiquePop/img/2026_07/Model_influence_Q/sigma_p_alagnon.png}%{a_juv.png}
     \caption{sigma\_p\_alagnon}
     \label{sigma_p_alagnon}
\end{figure}

\input{D:/Documents/Workspace_eclipse/ModeleDynamiquePop/tab/2026_07/Model_influence_Q/sigma_p_alagnon.tex}

sigma\_p\_alagnon représente l'écart-type autour du paramètre de blocage de la migration pour le secteur Alagnon.

\clearpage

%=================================================================================
\subsection{sigma\_p\_langeac}
%=================================================================================

\begin{figure}[ht]
     \includegraphics{D:/Documents/Workspace_eclipse/ModeleDynamiquePop/img/2026_07/Model_influence_Q/sigma_p_langeac.png}%{a_juv.png}
     \caption{sigma\_p\_langeac}
     \label{sigma_p_langeac}
\end{figure}

\input{D:/Documents/Workspace_eclipse/ModeleDynamiquePop/tab/2026_07/Model_influence_Q/sigma_p_langeac.tex}

sigma\_p\_langeac représente l'écart-type autour du paramètre de blocage de la migration pour le secteur Langeac-Poutès.

\clearpage

%=================================================================================
\subsection{sigma\_p\_poutes}
%=================================================================================

\begin{figure}[ht]
     \includegraphics{D:/Documents/Workspace_eclipse/ModeleDynamiquePop/img/2026_07/Model_influence_Q/sigma_p_poutes.png}%{a_juv.png}
     \caption{sigma\_p\_poutes}
     \label{sigma_p_poutes}
\end{figure}

\input{D:/Documents/Workspace_eclipse/ModeleDynamiquePop/tab/2026_07/Model_influence_Q/sigma_p_poutes.tex}

sigma\_p\_poutes représente l'écart-type autour du paramètre de blocage de la migration pour le secteur amont de Poutès.

\clearpage




%=================================================================================
\subsection{hel\_effect}
%=================================================================================

\begin{figure}[ht]
     \includegraphics{D:/Documents/Workspace_eclipse/ModeleDynamiquePop/img/2026_07/Model_influence_Q/hel_effect.png}%{a_juv.png}
     \caption{hel\_effect}
     \label{hel_effect}
\end{figure}

\input{D:/Documents/Workspace_eclipse/ModeleDynamiquePop/tab/2026_07/Model_influence_Q/hel_effect.tex}

hel\_effect représente l'effet de la méthodologie de comptage des nids (comptage à pied ou en bateau \textit{versus} comptage en hélicoptère).

\clearpage

%=================================================================================
\subsection{zone\_effect}
%================================================================================

zone\_effect représente l'effet secteur dans la relation nombre de nids par géniteur.
Les secteurs Vichy-Langeac / Alagnon / Langeac-Poutès se distinguent du secteur amont Poutès.


%-------------------------------------------
\subsubsection{zone\_effect\_Vichy}
%-------------------------------------------

\begin{figure}[ht]
     \includegraphics{D:/Documents/Workspace_eclipse/ModeleDynamiquePop/img/2026_07/Model_influence_Q/zone_effect_V.png}%{a_juv.png}
     \caption{zone\_effect\_V}
     \label{zone_effect_V}
\end{figure}
\clearpage


\begin{figure}[ht]
     \includegraphics{D:/Documents/Workspace_eclipse/ModeleDynamiquePop/img/2026_07/Model_influence_Q/zone_effect_A.png}%{a_juv.png}
     \caption{zone\_effect\_A}
     \label{zone_effect_A}
\end{figure}
\clearpage


%-------------------------------------------
\subsubsection{zone\_effect\_Langeac}
%-------------------------------------------

\begin{figure}[ht]
     \includegraphics{D:/Documents/Workspace_eclipse/ModeleDynamiquePop/img/2026_07/Model_influence_Q/zone_effect_L.png}%{a_juv.png}
     \caption{zone\_effect\_L}
     \label{zone_effect_L}
\end{figure}
\clearpage
%-------------------------------------------
\subsubsection{zone\_effect\_Poutes}
%-------------------------------------------

\begin{figure}[ht]
     \includegraphics{D:/Documents/Workspace_eclipse/ModeleDynamiquePop/img/2026_07/Model_influence_Q/zone_effect_P.png}%{a_juv.png}
     \caption{zone\_effect\_P}
     \label{zone_effect_P}
\end{figure}

\clearpage



%====================================
%\section{p\_comptage}
%====================================

% \begin{figure}[ht]
%      \includegraphics{D:/Documents/Workspace_eclipse/ModeleDynamiquePop/img/2026_07/Model_influence_Q/p_comptage.png}%{a_juv.png}
%      \caption{p\_comptage}
%      \label{p_comptage}
% \end{figure}
%
% \clearpage

%====================================
%\section{sigma\_p\_comptage}
%====================================

% \begin{figure}[ht]
%      \includegraphics{D:/Documents/Workspace_eclipse/ModeleDynamiquePop/img/2026_07/Model_influence_Q/sigma_p_comptage.png}%{a_juv.png}
%      \caption{sigma\_p\_comptage}
%      \label{sigma_p_comptage}
% \end{figure}
%
% \input{D:/Documents/Workspace_eclipse/ModeleDynamiquePop/tab/2026_07/Model_influence_Q/sigma_p_comptage.tex}
% \clearpage

%====================================
%\section{N\_Langeac_count}
%====================================


% \begin{figure}[ht]
%      \includegraphics{D:/Documents/Workspace_eclipse/ModeleDynamiquePop/img/2026_07/Model_influence_Q/N_langeac_count.png}%{a_juv.png}
%      \caption{N\_Langeac\_count : estimation de ce qu'on aurait compté si la station de Langeac avait été en service.}
%      \label{N_Langeac_count}
% \end{figure}
%
% \clearpage

%=================================================================
\subsection{Influence des débits sur la répartition des adultes}
%=================================================================

%=================================================================================
\subsubsection{beta\_Q\_L\_mont\_P}
%=================================================================================

\begin{figure}[ht]
     \includegraphics{D:/Documents/Workspace_eclipse/ModeleDynamiquePop/img/2026_07/Model_influence_Q/beta_Q_L_mont_P.png}
     \caption{beta\_Q\_L\_mont\_P}
     \label{beta_Q_L_mont_P}
\end{figure}

\input{D:/Documents/Workspace_eclipse/ModeleDynamiquePop/tab/2026_07/Model_influence_Q/beta_Q_L_mont_P.tex}

beta\_Q\_L\_mont\_P représente le coefficient modélisant l'influence des débits à Lempdes au printemps sur la répartition des adultes dans l'Alagnon.

\clearpage

%=================================================================================
\subsubsection{beta\_Q\_L\_mont\_A}
%=================================================================================


\begin{figure}[ht]
     \includegraphics{D:/Documents/Workspace_eclipse/ModeleDynamiquePop/img/2026_07/Model_influence_Q/beta_Q_L_mont_A.png}
     \caption{beta\_Q\_L\_mont\_A}
     \label{beta_Q_L_mont_A}
\end{figure}

\input{D:/Documents/Workspace_eclipse/ModeleDynamiquePop/tab/2026_07/Model_influence_Q/beta_Q_L_mont_A.tex}

beta\_Q\_L\_mont\_A représente le coefficient modélisant l'influence des débits à Lempdes à l'automne sur la répartition des adultes dans l'Alagnon.

\clearpage


%=================================================================================
\subsubsection{beta\_Q\_VB\_mont\_P}
%=================================================================================

\begin{figure}[ht]
     \includegraphics{D:/Documents/Workspace_eclipse/ModeleDynamiquePop/img/2026_07/Model_influence_Q/beta_Q_VB_mont_P.png}
     \caption{beta\_Q\_VB\_mont\_P}
     \label{beta_Q_VB_mont_P}
\end{figure}

\input{D:/Documents/Workspace_eclipse/ModeleDynamiquePop/tab/2026_07/Model_influence_Q/beta_Q_VB_mont_P.tex}

beta\_Q\_VB\_mont\_P représente le coefficient modélisant l'influence des débits à Vieille-Bridoude au printemps sur la répartition des adultes sur le secteur Langeac-Poutès.

\clearpage

%=================================================================================
\subsubsection{beta\_Q\_VB\_mont\_A}
%=================================================================================


\begin{figure}[ht]
     \includegraphics{D:/Documents/Workspace_eclipse/ModeleDynamiquePop/img/2026_07/Model_influence_Q/beta_Q_VB_mont_A.png}
     \caption{beta\_Q\_VB\_mont\_A}
     \label{beta_Q_VB_mont_A}
\end{figure}

\input{D:/Documents/Workspace_eclipse/ModeleDynamiquePop/tab/2026_07/Model_influence_Q/beta_Q_VB_mont_A.tex}

beta\_Q\_VB\_mont\_A représente le coefficient modélisant l'influence des débits à Vieille-Bridoude à l'automne sur la répartition des adultes sur le secteur Langeac-Poutès.

\clearpage


%=================================================================================
\subsubsection{beta\_Q\_P\_mont\_P}
%=================================================================================

\begin{figure}[ht]
     \includegraphics{D:/Documents/Workspace_eclipse/ModeleDynamiquePop/img/2026_07/Model_influence_Q/beta_Q_P_mont_P.png}
     \caption{beta\_Q\_P\_mont\_P}
     \label{beta_Q_P_mont_P}
\end{figure}

\input{D:/Documents/Workspace_eclipse/ModeleDynamiquePop/tab/2026_07/Model_influence_Q/beta_Q_P_mont_P.tex}

beta\_Q\_P\_mont\_P représente le coefficient modélisant l'influence des débits à Prades au printemps sur la répartition des adultes en amont Poutès

\clearpage

%=================================================================================
\subsubsection{beta\_Q\_P\_mont\_A}
%=================================================================================


\begin{figure}[ht]
     \includegraphics{D:/Documents/Workspace_eclipse/ModeleDynamiquePop/img/2026_07/Model_influence_Q/beta_Q_P_mont_A.png}
     \caption{beta\_Q\_P\_mont\_A}
     \label{beta_Q_P_mont_A}
\end{figure}

\input{D:/Documents/Workspace_eclipse/ModeleDynamiquePop/tab/2026_07/Model_influence_Q/beta_Q_P_mont_A.tex}

beta\_Q\_P\_mont\_A représente le coefficient modélisant l'influence des débits à Prades à l'automne sur la répartition des adultes en amont Poutès

\clearpage

\subsubsection{Représentation des séries de débits utilisées dans le modèle}



\begin{figure}[ht]
     \includegraphics{D:/Documents/Workspace_eclipse/ModeleDynamiquePop/img/2026_07/Model_influence_Q/plot_Q.png}
     \caption{représentation des séries de débits utilisées dans le modèle}
     \label{serie_Q}
\end{figure}

\clearpage

\subsubsection{Corrélogramme des séries de débits utilisées dans le modèle}

\begin{figure}[ht]
     \includegraphics{D:/Documents/Workspace_eclipse/ModeleDynamiquePop/img/2026_07/Model_influence_Q/corr_Q_rep.png}%{a_juv.png}
     \caption{Débits utilisés dans le modèle linéaire expliquant la
     répartition des adultes par les débits (Lempdes sur Alagnon pour les
     adultes migrant dans l'Alagnon, Vieille-Brioude pour les adultes migrant
     entre Langeac et Poutès et Prades pour les adultes migrant en amont de
     Poutès)}
     \label{corr_Q_rep}
\end{figure}


\clearpage

\section{Production des juvéniles (de l'adulte au juvénile)}

%=================================================================================
\subsection{a}
%=================================================================================

\begin{figure}[ht]
     \includegraphics{D:/Documents/Workspace_eclipse/ModeleDynamiquePop/img/2026_07/Model_influence_Q/a.png}%{a_juv.png}
     \caption{a}
     \label{a}
\end{figure}

\input{D:/Documents/Workspace_eclipse/ModeleDynamiquePop/tab/2026_07/Model_influence_Q/a.tex}
a représente le paramètre de pente de la relation Stock-Recrutement pour le compartiment sauvage, il représente donc le nombre d'oeufs produits par géniteur aux faibles déposes d'oeufs.
\clearpage

%=================================================================================
\subsection{a\_juv}
%=================================================================================

\begin{figure}[ht]
     \includegraphics{D:/Documents/Workspace_eclipse/ModeleDynamiquePop/img/2026_07/Model_influence_Q/a_juv.png}%{a_juv.png}
     \caption{a\_juv}
     \label{a_juv}
\end{figure}

\input{D:/Documents/Workspace_eclipse/ModeleDynamiquePop/tab/2026_07/Model_influence_Q/a_juv.tex}
a\_juv représente le paramètre de pente de la relation Stock-Recrutement pour le compartiment élevage. Il représente la survie de l'alevin déversé de juin à septembre.

\clearpage

%=================================================================================
\subsection{Rmax}
%=================================================================================

\begin{figure}[ht]
     \includegraphics{D:/Documents/Workspace_eclipse/ModeleDynamiquePop/img/2026_07/Model_influence_Q/Rmax.png}%{a_juv.png}
     \caption{Rmax}
     \label{Rmax}
\end{figure}

\input{D:/Documents/Workspace_eclipse/ModeleDynamiquePop/tab/2026_07/Model_influence_Q/Rmax.tex}
Rmax représente la capacité d'accueil du milieux. Plus précisément, il représente le recrutement en juvéniles sauvages qu'on peut espérer de mieux en médiane.
\clearpage

%=================================================================================
\subsection{p\_Rmax}
%=================================================================================
p\_Rmax représente la proportion d'habitats restante après implantation des juvéniles sauvages. Il intervient dans le mécanisme d'interaction réciproque entre juvéniles sauvages et juvéniles déversés. Le paramètre est décliné à l'échelle de chacun des secteurs du modèle.
%------------------------------------
\subsubsection{p\_Rmax\_V}
%------------------------------------

\begin{figure}[ht]
     \includegraphics{D:/Documents/Workspace_eclipse/ModeleDynamiquePop/img/2026_07/Model_influence_Q/p_Rmax_V.png}
     \caption{p\_Rmax\_V}
     \label{p_Rmax_V}
\end{figure}

\clearpage

%------------------------------------
\subsubsection{p\_Rmax\_A}
%------------------------------------

\begin{figure}[ht]
     \includegraphics{D:/Documents/Workspace_eclipse/ModeleDynamiquePop/img/2026_07/Model_influence_Q/p_Rmax_A.png}
     \caption{p\_Rmax\_A}
     \label{p_Rmax_A}
\end{figure}
\clearpage

%------------------------------------
\subsubsection{p\_Rmax\_L}
%------------------------------------

\begin{figure}[ht]
     \includegraphics{D:/Documents/Workspace_eclipse/ModeleDynamiquePop/img/2026_07/Model_influence_Q/p_Rmax_L.png}
     \caption{p\_Rmax\_L}
     \label{p_Rmax_L}
\end{figure}
\clearpage

%------------------------------------
\subsubsection{p\_Rmax\_P}
%------------------------------------

\begin{figure}[ht]
     \includegraphics{D:/Documents/Workspace_eclipse/ModeleDynamiquePop/img/2026_07/Model_influence_Q/p_Rmax_P.png}
     \caption{p\_Rmax\_P}
     \label{p_Rmax_P}
\end{figure}
\clearpage

%=================================================================================
%{mu\_tau}
%=================================================================================

%\begin{figure}[ht]
%     \includegraphics{D:/Documents/Workspace_eclipse/ModeleDynamiquePop/img/2026_07/Model_influence_Q/mu_tau.png}%{a_juv.png}
%     \caption{mu\_tau}
%     \label{mu_tau}
%\end{figure}

%\input{D:/Documents/Workspace_eclipse/ModeleDynamiquePop/tab/2026_07/Model_influence_Q/mu_tau.tex}

%mu\_tau est un hyperparamètre intervenant dans la précision des densités de
% juvéniles.

%\clearpage

%=================================================================================
\subsection{alpha\_sigma}
%=================================================================================

\begin{figure}[ht]
     \includegraphics{D:/Documents/Workspace_eclipse/ModeleDynamiquePop/img/2026_07/Model_influence_Q/alpha_sigma.png}%{a_juv.png}
     \caption{alpha\_sigma}
     \label{alpha_sigma}
\end{figure}

\input{D:/Documents/Workspace_eclipse/ModeleDynamiquePop/tab/2026_07/Model_influence_Q/alpha_sigma.tex}

alpha\_sigma est un hyperparamètre intervenant dans l'écart-type des densités de
juvéniles, des passages à Vichy et des probabilité de migrer dans les
différentes zones. Toutes ces variances ont une même structure hiérarchique.

\clearpage

%=================================================================================
\subsection{mu\_sigma}
%=================================================================================

\begin{figure}[ht]
     \includegraphics{D:/Documents/Workspace_eclipse/ModeleDynamiquePop/img/2026_07/Model_influence_Q/mu_sigma.png}%{a_juv.png}
     \caption{mu\_sigma}
     \label{mu_sigma}
\end{figure}

\input{D:/Documents/Workspace_eclipse/ModeleDynamiquePop/tab/2026_07/Model_influence_Q/mu_sigma.tex}

mu\_sigma est un hyperparamètre intervenant dans la précision des densités de
juvéniles.

\clearpage

%=================================================================================
%{beta\_tau}
%=================================================================================

%\begin{figure}[ht]
%     \includegraphics{D:/Documents/Workspace_eclipse/ModeleDynamiquePop/img/2026_07/Model_influence_Q/beta_tau.png}%{a_juv.png}
%     \caption{beta\_tau}
%     \label{beta_tau}
%\end{figure}
%
%\input{D:/Documents/Workspace_eclipse/ModeleDynamiquePop/tab/2026_07/Model_influence_Q/beta_tau.tex}
%
%beta\_tau est un hyperparamètre intervenant dans la précision des densités de juvéniles.
%
%\clearpage

%=================================================================================
%\section{rho\_D}
%=================================================================================
%\subsection{rho\_D}

%\begin{figure}[ht]
%\includegraphics{D:/Documents/Workspace_eclipse/ModeleDynamiquePop/img/2026_07/Model_influence_Q/rho_D.png}%{a_juv.png}
%		\caption{rho\_D}
%\label{rho_D}
%\end{figure}
%
%
%\input{D:/Documents/Workspace_eclipse/ModeleDynamiquePop/tab/2026_07/Model_influence_Q/rho_D.tex}
%
%rho\_D représente le coefficient de corrélation entre les variances du recrutement sauvage et élevage.
%\clearpage

%=================================================================================
%\section{nu\_wild}
%=================================================================================

%\begin{figure}[ht]
%     \includegraphics{D:/Documents/Workspace_eclipse/ModeleDynamiquePop/img/2026_07/Model_influence_Q/nu_wild.png}
%     \caption{nu\_wild}
%     \label{nu_wild}
%\end{figure}
%
%\input{D:/Documents/Workspace_eclipse/ModeleDynamiquePop/tab/2026_07/Model_influence_Q/nu_wild.tex}
%\clearpage


%=================================================================================
\subsection{nu\_d}
%=================================================================================

\begin{figure}[ht]
     \includegraphics{D:/Documents/Workspace_eclipse/ModeleDynamiquePop/img/2026_07/Model_influence_Q/nu_d.png}
     \caption{nu\_d}
     \label{nu_d}
\end{figure}

\input{D:/Documents/Workspace_eclipse/ModeleDynamiquePop/tab/2026_07/Model_influence_Q/nu_d.tex}

nu\_d représente le différentiel de production en juvéniles des différents secteurs du modèle. Il est utilisé dans la relation Stock-Recrutement (Beverton et Holt).
\clearpage

%=================================================================================
\subsection{sigma\_wild\_moy}
%=================================================================================

\begin{figure}[ht]
     \includegraphics{D:/Documents/Workspace_eclipse/ModeleDynamiquePop/img/2026_07/Model_influence_Q/sigma_wild_moy.png}%{a_juv.png}
     \caption{sigma\_wild\_moy}
     \label{sigma_wild_moy}
\end{figure}


\input{D:/Documents/Workspace_eclipse/ModeleDynamiquePop/tab/2026_07/Model_influence_Q/sigma_wild_moy.tex}

sigma\_wild\_moy représente l'écart-type des densité de juvéniles sauvages à l'échelle des secteurs. Ce paramètre est utilisé dans la matrice de variance-covariance qui régit l'interaction réciproque entre les juvéniles sauvages et les juvéniles déversés.
\clearpage

%=================================================================================
\subsection{sigma\_juv\_moy}
%=================================================================================

\begin{figure}[ht]
     \includegraphics{D:/Documents/Workspace_eclipse/ModeleDynamiquePop/img/2026_07/Model_influence_Q/sigma_juv_moy.png}%{a_juv.png}
     \caption{sigma\_juv\_moy}
     \label{sigma_juv_moy}
\end{figure}


\input{D:/Documents/Workspace_eclipse/ModeleDynamiquePop/tab/2026_07/Model_influence_Q/sigma_juv_moy.tex}

sigma\_juv\_moy représente l'écart-type des densité de juvéniles déversés à l'échelle des secteurs. Ce paramètre est utilisé dans la matrice de variance-covariance qui régit l'interaction réciproque entre les juvéniles sauvages et les juvéniles déversés.
\clearpage


%=================================================================================
\subsection{sigma\_egg\_moy}
%=================================================================================

\begin{figure}[ht]
     \includegraphics{D:/Documents/Workspace_eclipse/ModeleDynamiquePop/img/2026_07/Model_influence_Q/sigma_egg_moy.png}%{a_juv.png}
     \caption{sigma\_egg\_moy}
     \label{sigma_egg_moy}
\end{figure}

\input{D:/Documents/Workspace_eclipse/ModeleDynamiquePop/tab/2026_07/Model_influence_Q/sigma_egg_moy.tex}

sigma\_egg\_moy représente l'écart-type des densités de juvéniles issus des incubateurs de terrain à l'échelle des secteurs.

\clearpage

%=================================================================================
\subsection{sigma\_wild\_site}
%=================================================================================

\begin{figure}[ht]
     \includegraphics{D:/Documents/Workspace_eclipse/ModeleDynamiquePop/img/2026_07/Model_influence_Q/sigma_wild_site.png}%{a_juv.png}
     \caption{sigma\_wild\_site}
     \label{sigma_wild_site}
\end{figure}

\input{D:/Documents/Workspace_eclipse/ModeleDynamiquePop/tab/2026_07/Model_influence_Q/sigma_wild_site.tex}

sigma\_wild\_site représente l'écart-type des densité de juvéniles sauvages à l'échelle des sites (stations de pêche électrique).

\clearpage



%=================================================================================
\subsection{sigma\_juv\_site}
%=================================================================================

\begin{figure}[ht]
     \includegraphics{D:/Documents/Workspace_eclipse/ModeleDynamiquePop/img/2026_07/Model_influence_Q/sigma_juv_site.png}%{a_juv.png}
     \caption{sigma\_juv\_site}
     \label{sigma_juv_site}
\end{figure}

\input{D:/Documents/Workspace_eclipse/ModeleDynamiquePop/tab/2026_07/Model_influence_Q/sigma_juv_site.tex}

sigma\_juv\_site représente l'écart-type des densité de juvéniles déversés à l'échelle des sites (stations de pêche électrique).
\clearpage


%=================================================================================
\subsection{sigma\_egg\_site}
%=================================================================================

\begin{figure}[ht]
     \includegraphics{D:/Documents/Workspace_eclipse/ModeleDynamiquePop/img/2026_07/Model_influence_Q/sigma_egg_site.png}%{a_juv.png}
     \caption{sigma\_egg\_site}
     \label{sigma_egg_site}
\end{figure}

\input{D:/Documents/Workspace_eclipse/ModeleDynamiquePop/tab/2026_07/Model_influence_Q/sigma_egg_site.tex}

sigma\_egg\_site représente l'écart-type des densités de juvéniles issus des incubateurs de terrain à l'échelle des sites (stations de pêche électrique).

\clearpage


%====================================
\subsection{d\_wild\_moy}
%====================================

d\_wild\_moy représente la densité de juvénile sauvage à l'échelle des secteurs. Il est décliné pour chacun des 4 secteurs.

%------------------------------------
\subsubsection{d\_wild\_moy\_Vichy}
%------------------------------------

\begin{figure}[ht]
     \includegraphics{D:/Documents/Workspace_eclipse/ModeleDynamiquePop/img/2026_07/Model_influence_Q/d_wild_moy_V.png}%{a_juv.png}
     \caption{d\_wild\_moy\_V}
     \label{d_wild_moy_V}
\end{figure}

\clearpage


%------------------------------------
\subsubsection{d\_wild\_moy\_Alagnon}
%------------------------------------

\begin{figure}[ht]
     \includegraphics{D:/Documents/Workspace_eclipse/ModeleDynamiquePop/img/2026_07/Model_influence_Q/d_wild_moy_A.png}%{a_juv.png}
     \caption{d\_wild\_moy\_A}
     \label{d_wild_moy_A}
\end{figure}

\clearpage

%------------------------------------
\subsubsection{d\_wild\_moy\_Langeac}
%------------------------------------

\begin{figure}[ht]
     \includegraphics{D:/Documents/Workspace_eclipse/ModeleDynamiquePop/img/2026_07/Model_influence_Q/d_wild_moy_L.png}%{a_juv.png}
     \caption{d\_wild\_moy\_L}
     \label{d_wild_moy_L}
\end{figure}

\clearpage

%------------------------------------
\subsubsection{d\_wild\_moy\_Poutes}
%------------------------------------

\begin{figure}[ht]
     \includegraphics{D:/Documents/Workspace_eclipse/ModeleDynamiquePop/img/2026_07/Model_influence_Q/d_wild_moy_P.png}%{a_juv.png}
     \caption{d\_wild\_moy\_P}
     \label{d_wild_moy_P}
\end{figure}

\clearpage


%====================================
\subsection{d\_juv\_moy}
%====================================

d\_juv\_moy représente la densité de juvéniles déversés à l'échelle des secteurs. Il est décliné à l'échelle des 4 secteurs du modèle.

%------------------------------------
\subsubsection{d\_juv\_moy\_Vichy}
%------------------------------------

\begin{figure}[ht]
     \includegraphics{D:/Documents/Workspace_eclipse/ModeleDynamiquePop/img/2026_07/Model_influence_Q/d_juv_moy_V.png}%{a_juv.png}
     \caption{d\_juv\_moy\_V}
     \label{d_juv_moy_V}
\end{figure}

\clearpage

%------------------------------------
\subsubsection{d\_juv\_moy\_Alagnon}
%------------------------------------

\begin{figure}[ht]
     \includegraphics{D:/Documents/Workspace_eclipse/ModeleDynamiquePop/img/2026_07/Model_influence_Q/d_juv_moy_A.png}%{a_juv.png}
     \caption{d\_juv\_moy\_A}
     \label{d_juv_moy_A}
\end{figure}

\clearpage

%------------------------------------
\subsubsection{d\_juv\_moy\_Langeac}
%------------------------------------

\begin{figure}[ht]
     \includegraphics{D:/Documents/Workspace_eclipse/ModeleDynamiquePop/img/2026_07/Model_influence_Q/d_juv_moy_L.png}%{a_juv.png}
     \caption{d\_juv\_moy\_L}
     \label{d_juv_moy_L}
\end{figure}

\clearpage

%------------------------------------
\subsubsection{d\_juv\_moy\_Poutes}
%------------------------------------

\begin{figure}[ht]
     \includegraphics{D:/Documents/Workspace_eclipse/ModeleDynamiquePop/img/2026_07/Model_influence_Q/d_juv_moy_P.png}%{a_juv.png}
     \caption{d\_juv\_moy\_P}
     \label{d_juv_moy_P}
\end{figure}

\clearpage

%====================================
\subsection{d\_egg\_moy}
%====================================

d\_egg\_moy représente la densité de juvéniles issus des incubateurs de terrain à l'échelle des secteurs. Il est décliné sur les 2 secteurs ayant abrité dans le passé ou le présent un ou des incubateurs de terrain.

%------------------------------------
\subsubsection{d\_egg\_moy\_Vichy}
%------------------------------------

\begin{figure}[ht]
     \includegraphics{D:/Documents/Workspace_eclipse/ModeleDynamiquePop/img/2026_07/Model_influence_Q/d_egg_moy_V.png}%{a_juv.png}
     \caption{d\_egg\_moy\_V}
     \label{d_egg_moy_V}
\end{figure}

\clearpage

%------------------------------------
\subsubsection{d\_egg\_moy\_Langeac}
%------------------------------------

\begin{figure}[ht]
     \includegraphics{D:/Documents/Workspace_eclipse/ModeleDynamiquePop/img/2026_07/Model_influence_Q/d_egg_moy_L.png}%{a_juv.png}
     \caption{d\_egg\_moy\_L}
     \label{d_egg_moy_L}
\end{figure}

\clearpage

%====================================
\subsection{res\_wild\_moy}
%====================================

res\_wild\_moy représente les résidus de la densité de juvénile sauvage à l'échelle de chacun des secteurs du modèle.

%------------------------------------
\subsubsection{res\_wild\_moy\_Vichy\_standardisé}
%------------------------------------

\begin{figure}[ht]
     \includegraphics{D:/Documents/Workspace_eclipse/ModeleDynamiquePop/img/2026_07/Model_influence_Q/res_wild_moy_V_std.png}
     \caption{res\_wild\_moy\_V\_std}
     \label{res_wild_moy_V_std}
\end{figure}

\clearpage


%------------------------------------
\subsubsection{res\_wild\_moy\_Alagnon\_standardisé}
%------------------------------------

\begin{figure}[ht]
     \includegraphics{D:/Documents/Workspace_eclipse/ModeleDynamiquePop/img/2026_07/Model_influence_Q/res_wild_moy_A_std.png}%{a_juv.png}
     \caption{res\_wild\_moy\_A\_std}
     \label{res_wild_moy_A_std}
\end{figure}

\clearpage


%------------------------------------
\subsubsection{res\_wild\_moy\_Langeac\_standardisé}
%------------------------------------

\begin{figure}[ht]
     \includegraphics{D:/Documents/Workspace_eclipse/ModeleDynamiquePop/img/2026_07/Model_influence_Q/res_wild_moy_L_std.png}%{a_juv.png}
     \caption{res\_wild\_moy\_L\_std}
     \label{res_wild_moy_L_std}
\end{figure}

\clearpage


%------------------------------------
\subsubsection{res\_wild\_moy\_Poutes\_standardisé}
%------------------------------------

\begin{figure}[ht]
     \includegraphics{D:/Documents/Workspace_eclipse/ModeleDynamiquePop/img/2026_07/Model_influence_Q/res_wild_moy_P_std.png}%{a_juv.png}
     \caption{res\_wild\_moy\_P\_std}
     \label{res_wild_moy_P_std}
\end{figure}

\clearpage


%====================================
\subsection{res\_juv\_moy}
%====================================

res\_juv\_moy représente les résidus de la densité de juvénile déversés à l'échelle de chacun des secteurs du modèle.

%------------------------------------
\subsubsection{res\_juv\_moy\_Vichy\_standardisé}
%------------------------------------

\begin{figure}[ht]
     \includegraphics{D:/Documents/Workspace_eclipse/ModeleDynamiquePop/img/2026_07/Model_influence_Q/res_juv_moy_V_std.png}%{a_juv.png}
     \caption{res\_juv\_moy\_V\_std}
     \label{res_juv_moy_V_std}
\end{figure}

\clearpage


%------------------------------------
\subsubsection{res\_juv\_moy\_Alagnon\_standardisé}
%------------------------------------

\begin{figure}[ht]
     \includegraphics{D:/Documents/Workspace_eclipse/ModeleDynamiquePop/img/2026_07/Model_influence_Q/res_juv_moy_A_std.png}%{a_juv.png}
     \caption{res\_juv\_moy\_A\_std}
     \label{res_juv_moy_A_std}
\end{figure}

\clearpage



%------------------------------------
\subsubsection{res\_juv\_moy\_Langeac\_standardisé}
%------------------------------------

\begin{figure}[ht]
     \includegraphics{D:/Documents/Workspace_eclipse/ModeleDynamiquePop/img/2026_07/Model_influence_Q/res_juv_moy_L_std.png}%{a_juv.png}
     \caption{res\_juv\_moy\_L\_std}
     \label{res_juv_moy_L_std}
\end{figure}

\clearpage

%------------------------------------
\subsubsection{res\_juv\_moy\_Poutes\_standardisé}
%------------------------------------

\begin{figure}[ht]
     \includegraphics{D:/Documents/Workspace_eclipse/ModeleDynamiquePop/img/2026_07/Model_influence_Q/res_juv_moy_P_std.png}%{a_juv.png}
     \caption{res\_juv\_moy\_P\_std}
     \label{res_juv_moy_P_std}
\end{figure}

\clearpage


%=================================================================================
\subsection{Corrélations spatiales des résidus de densité de juvéniles sauvages}
%=================================================================================


\begin{figure}[ht]
     \includegraphics{D:/Documents/Workspace_eclipse/ModeleDynamiquePop/img/2026_07/Model_influence_Q/corr_res_wild.png}
     \caption{corr\_res\_wild}
     \label{corr_res_wild}
\end{figure}

\clearpage

%=================================================================================
\subsection{Corrélations spatiales des résidus de densité de juvéniles élevages}
%=================================================================================

\begin{figure}[ht]
     \includegraphics{D:/Documents/Workspace_eclipse/ModeleDynamiquePop/img/2026_07/Model_influence_Q/corr_res_juv.png}
     \caption{corr\_res\_juv}
     \label{corr_res_juv}
\end{figure}

\clearpage


%=================================================================================
%\section{level\_s}
%=================================================================================

% \begin{figure}[ht]
%     \includegraphics{D:/Documents/Workspace_eclipse/ModeleDynamiquePop/img/2026_07/Model_influence_Q/level_s.png}%{a_juv.png}
%     \caption{level\_s}
%     \label{level_s}
% \end{figure}
%
% \input{D:/Documents/Workspace_eclipse/ModeleDynamiquePop/tab/2026_07/Model_influence_Q/level_s.tex}
% \clearpage


%=================================================================================
%\section{tau\_Q}
%=================================================================================

% \begin{figure}[ht]
%      \includegraphics{D:/Documents/Workspace_eclipse/ModeleDynamiquePop/img/2026_07/Model_influence_Q/tau_Q.png}
%      \caption{tau\_Q :précision de la loi normale dans laquelle on tire L_s_juv2ad}
%      \label{tau_Q}
% \end{figure}
%
% \input{D:/Documents/Workspace_eclipse/ModeleDynamiquePop/tab/2026_07/Model_influence_Q/tau_Q.tex}
%\clearpage




%=================================================================================
%\section{rho\_poutes}
%=================================================================================

%\begin{figure}[ht]
%     \includegraphics{D:/Documents/Workspace_eclipse/ModeleDynamiquePop/img/2026_07/Model_influence_Q/rho_poutes.png}%{a_juv.png}
%     \caption{rho\_poutes}
%     \label{rho_poutes}
%\end{figure}
%
%\input{D:/Documents/Workspace_eclipse/ModeleDynamiquePop/tab/2026_07/Model_influence_Q/rho_poutes.tex}
%\clearpage







\end{document}
